% Rédaction : BANDE Hamidou ; relecture : SOMBIE Roseline
\chapter{Données et base PostgreSQL}\label{chap:donnees}

\section{Des données fictives, piégées exprès}
Le générateur produit 1000 transactions et ajoute 10 doublons de référence, soit 1010 lignes.
\begin{itemize}
  \item programme : \texttt{data/scripts/generate-transactions.js} ;
  \item fichier produit : \texttt{data/raw/transactions\_brutes.csv} ;
  \item liste des anomalies injectées : \texttt{data/raw/anomalies\_injectees.json}.
\end{itemize}
Il utilise une graine fixe (2026) : relancé sur n'importe quel poste, il redonne exactement le même
fichier. Le résultat attendu de l'ETL est donc connu à la ligne près.

Une chaîne de contrôle qui ne reçoit que des données propres ne prouve rien. Le générateur injecte
donc deux sortes d'anomalies, chacune listée avec le motif attendu :
\begin{itemize}
  \item \textbf{54 anomalies bloquantes} : la donnée est absente, fausse ou ambiguë (montant négatif,
        devise \texttt{FCFA}, date impossible, référence mal formée\ldots). La ligne est rejetée avec un motif ;
  \item \textbf{26 anomalies corrigeables} : l'intention est certaine (minuscules, espaces, virgule
        décimale, séparateur de milliers, date au format français). La ligne est corrigée puis acceptée.
\end{itemize}

\begin{table}[H]
\centering\small
\caption{Composition du jeu de données}
\begin{tabular}{|l|r|}
\hline
\rowcolor{bluedark}\textcolor{white}{\textbf{Catégorie}} & \textcolor{white}{\textbf{Lignes}} \\
\hline
Lignes brutes & 1010 \\
\hline
\quad dont doublons de référence & 10 \\
\hline
\quad dont anomalies bloquantes & 54 \\
\hline
\quad dont anomalies corrigeables (acceptées après correction) & 26 \\
\hline
Lignes valides attendues & 946 \\
\hline
Lignes rejetées attendues (54 + 10) & 64 \\
\hline
\end{tabular}
\end{table}

Un ETL de référence écrit en JavaScript (\texttt{data/scripts/reference-etl.js}) applique les
mêmes règles que KNIME et produit le résultat attendu dans \texttt{data/expected/}. Il sert
d'oracle : le comparateur \texttt{npm run data:compare} confronte la sortie de KNIME à ce résultat,
ligne par ligne (chapitre~\ref{chap:knime}).

\section{Le schéma de la base}
Le script \texttt{db/init/01\_schema.sql} est exécuté automatiquement par PostgreSQL à son premier
démarrage. Il crée trois tables :
\begin{itemize}
  \item \texttt{transactions} : les transactions valides, venues de KNIME (\texttt{source = 'KNIME'}) ou de l'API (\texttt{source = 'API'}) ;
  \item \texttt{transactions\_rejetees} : les lignes rejetées, stockées en texte brut avec leur \texttt{motif\_rejet} ;
  \item \texttt{chargements} : le journal de chaque chargement (lignes lues, insérées, mises à jour, rejetées).
\end{itemize}

\capture{hamidou_01_schema.png}{Les trois tables créées au démarrage de PostgreSQL}{fig:schema}

La table \texttt{transactions} porte les règles métier sous forme de contraintes : format de la
référence (\texttt{\^{}TXN-[0-9]\{6\}\$}) et des comptes, valeurs autorisées pour le type, la devise,
le statut et le canal, \texttt{montant > 0}, compte destination obligatoire pour un virement et
différent du compte source. Six index accélèrent les filtres de l'API (date, type, statut, comptes,
montant).

Ces contraintes sont une \textbf{dernière barrière} : même un \texttt{INSERT} écrit à la main est refusé.

\capture{hamidou_02_insert_refuse.png}{Un montant négatif est refusé par la contrainte \texttt{transactions\_montant\_check}}{fig:insert-refuse}

\section{Trois barrières de qualité}
\begin{enumerate}
  \item \textbf{KNIME} applique toutes les règles et motive chaque rejet ;
  \item \textbf{le chargeur} revalide chaque ligne « valide » avec les règles de l'API ; une ligne
        refusée à ce stade est enregistrée avec le motif \texttt{REJET\_CHARGEMENT} (cela ne doit pas arriver) ;
  \item \textbf{PostgreSQL} refuse toute ligne qui viole une contrainte.
\end{enumerate}

\section{Un chargement idempotent et transactionnel}
Le chargeur \texttt{api/src/scripts/load-csv.js} lit les deux fichiers produits par KNIME et les
charge dans \textbf{une seule transaction SQL} : soit tout est chargé, soit rien. Il utilise
\texttt{INSERT \ldots{} ON CONFLICT (reference) DO UPDATE} : relancer le chargement met à jour les
lignes existantes au lieu de créer des doublons. Une transaction créée par l'API n'est jamais
écrasée par un rechargement KNIME.

\capture{hamidou_03_idempotence.png}{Deuxième chargement : aucune insertion, 946 mises à jour, 64 rejets}{fig:idempotence}

\begin{table}[H]
\centering\small
\caption{Idempotence du chargement (poste de BANDE, 28/10/2026)}
\begin{tabular}{|l|r|r|r|}
\hline
\rowcolor{bluedark}\textcolor{white}{\textbf{Chargement}} & \textcolor{white}{\textbf{Insérées}} & \textcolor{white}{\textbf{Mises à jour}} & \textcolor{white}{\textbf{Rejetées}} \\
\hline
1\textsuperscript{er} passage & 946 & 0 & 64 \\
\hline
2\textsuperscript{e} passage & 0 & 946 & 64 \\
\hline
\end{tabular}
\end{table}
